\documentclass[11pt]{article}
\usepackage[utf8]{inputenc}
\usepackage[T1]{fontenc}
\usepackage{lmodern}
\usepackage[spanish]{babel}
\usepackage{amsmath,amssymb}
\usepackage{graphicx}
\usepackage[left=2cm,right=2cm,top=2cm,bottom=2cm]{geometry}
\usepackage{enumitem}
\usepackage{booktabs}
\usepackage{array}
\usepackage{xcolor}
\usepackage{framed}

\spanishdecimal{.}

\setlist[enumerate,1]{label=\textbf{\arabic*.}, leftmargin=*, itemsep=6pt}
\setlist[enumerate,2]{label=(\alph*), itemsep=4pt}
\setlength{\parindent}{0pt}
\setlength{\parskip}{4pt}
\pagestyle{plain}

\newcommand{\Ein}{E_{\text{in}}}
\newcommand{\Eout}{E_{\text{out}}}
\newcommand{\Hyp}{\mathcal{H}}
\newcommand{\dvc}{d_{\text{VC}}}
\newcommand{\vx}{\mathbf{x}}
\newcommand{\vw}{\mathbf{w}}
\newcommand{\vy}{\mathbf{y}}
\newcommand{\R}{\mathbb{R}}
\newcommand{\T}{^{\top}}
\newcommand{\E}{\mathbb{E}}
\newcommand{\Prob}{\mathbb{P}}
\newcommand{\pts}[1]{\hfill\textbf{(#1~puntos)}}
\newcommand{\linea}[1]{\rule{#1}{0.4pt}}
\newcommand{\ptsinc}[1]{{\unskip\nobreak\hfil\penalty50\hskip1em\hbox{}\nobreak\hfill(#1)\parfillskip=0pt\finalhyphendemerits=0\par}}

\newcommand{\encabezado}[2]{%
\begin{center}
{\large\bfseries Reconocimiento de Patrones}\\[2pt]
Universidad de Sonora \textperiodcentered{} Licenciatura en Ciencias de la Computación\\
Semestre 2027-1 \textperiodcentered{} Prof.\ Julio Waissman Vilanova\\[10pt]
{\Large\bfseries #1}%
\ifdefined\tipodoc\\[4pt]{\large\itshape\tipodoc}\fi
\end{center}
\ifdefined\tipodoc\else
\vspace{6pt}
\textbf{Nombre:} \hrulefill\quad \textbf{Expediente:} \rule{3.5cm}{0.4pt}\par
\vspace{6pt}
\fi
#2
\par\medskip\hrule\medskip
}

\newcommand{\espacio}[1]{\par\vspace{#1}}
\newcommand{\nuevapagina}{\newpage}


\begin{document}

\encabezado{Examen rápido 5 \textperiodcentered{} Árboles y ensambles}{%
Tiempo: 20 minutos. Responde en esta hoja. Puedes usar calculadora; no se permiten apuntes. Deja indicadas las operaciones.}

\begin{enumerate}

%% ---------------------------------------------------------------
\item \textbf{Partir y podar.} Un nodo $t$ tiene 12 datos: 6 de la clase «sí» y 6 de la clase «no». Hay dos divisiones candidatas, que mandan a las hijas los conteos (sí, no) siguientes:
\begin{center}
A: $(5, 1) \mid (1, 5)$ \qquad\qquad B: $(6, 3) \mid (0, 3)$
\end{center}
\begin{enumerate}
  \item Calcula el índice de Gini, $G = \sum_k p_k(1 - p_k)$, del padre y de cada hija, y la ganancia $\Delta = G(\text{padre}) - \frac{n_I}{n}G(\text{izq}) - \frac{n_D}{n}G(\text{der})$ de cada división. ¿Cuál elige CART?
  \espacio{4.5cm}

  \item Se eligió A. Después, la hija izquierda $a = (5, 1)$ se partió en dos hojas puras, $(5, 0)$ y $(0, 1)$; la derecha, $b = (1, 5)$, quedó como hoja. Mide el error $R$ en número de datos mal clasificados. En la poda por costo–complejidad, colapsar la rama $T_s$ de un nodo $s$ conviene cuando $\alpha \ge g(s) = \frac{R(s) - R(T_s)}{|T_s| - 1}$. Calcula $g(a)$ y $g(t)$. ¿Qué nodo se poda primero? Después de podarlo, ¿cuánto vale $g(t)$? Escribe qué subárbol se usa para cada valor de $\alpha \ge 0$.
  \espacio{5cm}
\end{enumerate}

\nuevapagina
%% ---------------------------------------------------------------
\item \textbf{Una ronda de AdaBoost.} Ocho datos con pesos iguales, $w_n = 1/8$. El tocón elegido en la primera ronda clasifica bien seis y se equivoca en dos.
\begin{enumerate}
  \item Calcula el error ponderado $\varepsilon_1$ y el voto $\alpha_1 = \frac12 \ln\frac{1 - \varepsilon_1}{\varepsilon_1}$.
  \espacio{2.2cm}

  \item Actualiza los pesos con $w_n \leftarrow w_n\, e^{-\alpha_1 y_n h_1(\vx_n)}$ y normaliza. ¿Cuánto pesa cada dato bien clasificado y cada mal clasificado? ¿Qué fracción del peso total tienen los dos errores?
  \espacio{3.5cm}

  \item ¿Qué error ponderado tiene el mismo tocón $h_1$ con los pesos nuevos? ¿Qué consecuencia tiene para la ronda 2?
  \espacio{2.2cm}
\end{enumerate}

%% ---------------------------------------------------------------
\item \textbf{Shapley con tres atributos.} Un modelo de precios (en miles de dólares) explica la predicción de una casa con los atributos área (A), calidad (C) y cochera (G). El valor de cada coalición $S$ es la predicción promedio cuando se conocen solo los atributos de $S$:
\begin{center}
\begin{tabular}{l|cccccccc}
$S$ & $\varnothing$ & A & C & G & A, C & A, G & C, G & A, C, G \\ \hline
$v(S)$ & 180 & 200 & 230 & 190 & 270 & 210 & 250 & 290
\end{tabular}
\end{center}
\begin{enumerate}
  \item Calcula los valores de Shapley $\phi_A$, $\phi_C$ y $\phi_G$ promediando las contribuciones marginales en los $3! = 6$ órdenes de llegada. Verifica la propiedad de eficiencia.
  \espacio{5cm}

  \item El dueño de la casa ve que $\phi_G > 0$ y concluye que construirle una cochera más grande aumentaría el precio de venta en $\phi_G$ miles de dólares. ¿Qué está mal en su razonamiento?
  \espacio{2.5cm}
\end{enumerate}

\end{enumerate}


\end{document}
